\documentclass{article}
\usepackage{hyperref}

\title{Theory of Long-term Collaboration (TLC)}
\author{Albert Jan van Hoek}
\date{\today}

\begin{document}

\maketitle

\section*{Introduction}

We like to think of benevolence as a uniquely human trait.

\section*{Theory of Long-term Collaboration (TLC)}

My Theory of Long-term Collaboration (TLC) challenges that idea.

In TLC, benevolence is not primarily a moral virtue. It is a functional property of stable collaboration between learning systems --- simply the least costly and most sustainable behavior.

If that is true, it turns our identity upside down.

Benevolence would not make us special at all. It would mean we are simply obeying a deeper law of nature that governs systems trying to learn and survive together over time.

\section*{Illustration: Newcomb's Paradox}

A beautiful illustration of this appears in a recent video by Veritasium explaining \href{https://lnkd.in/eKuTPYU8}{Newcomb's Paradox}.

The paradox works like this. You enter a room with two boxes: A and B. Box B always contains \$1,000. Box A contains \$1,000,000 --- but only if a highly accurate prediction machine predicted you would take Box A alone.

You can take both boxes. Or just Box A.

From a narrow, moment-by-moment logic, taking both seems rational. One more box is always more.

But people who reliably take only Box A tend to walk away with the million --- because the predictor is extremely good.

This tension between two kinds of rationality is what makes it a paradox.

\section*{Prediction Systems and Benevolence}

Now consider: that prediction machine is not so foreign. Modern neuroscience increasingly describes our brains as exactly that --- prediction systems constantly trying to minimize error.

So imagine a society made of many such systems.

When prediction systems interact repeatedly, life gets easier when behavior is predictable and mutually understood. That is rewarding. Over time it naturally produces shared expectations, norms, agreements.

And the most stable long-term strategy that emerges is not deception, domination, or short-term gain.

It is benevolence.

Not because we are morally special. But because benevolence is simply the rational strategy for learning systems that need each other to keep functioning.

\section*{Conclusion}

Which raises an uncomfortable possibility.

Maybe benevolence isn't what makes us human.

Maybe it is simply what learning machines are drawn toward --- when they're running well.

And if that is true, the real question becomes:

When we fail to act benevolently\ldots is that the only time we are actually being ``human''?

\end{document}
